\documentclass[12pt]{article}
\usepackage{amsmath, amssymb}
\usepackage{geometry}
\usepackage{hyperref}
\usepackage{listings}
\usepackage{xcolor}
\usepackage{booktabs}
\usepackage{enumitem}

\geometry{margin=1in}

\definecolor{codegray}{rgb}{0.95,0.95,0.95}
\definecolor{codegreen}{rgb}{0.1,0.5,0.1}
\definecolor{codepurple}{rgb}{0.4,0.0,0.6}
\lstset{
  backgroundcolor=\color{codegray},
  basicstyle=\ttfamily\small,
  keywordstyle=\color{codepurple}\bfseries,
  commentstyle=\color{codegreen},
  breaklines=true,
  frame=single,
  framesep=4pt,
  language=Python
}

\title{NFW + Strong Lensing Pipeline:\\Diagnosis and Repair of Three Compounding Bugs}
\author{Research Notes}
\date{\today}

\begin{document}
\maketitle

\tableofcontents
\newpage

% ─────────────────────────────────────────────────────────────────
\section{Background and Problem Statement}
% ─────────────────────────────────────────────────────────────────

The gravitational lensing reconstruction pipeline (\texttt{fit\_lensing\_field}) supports two
lens models --- SIS and NFW --- and can optionally incorporate multiply-imaged strong-lensing
(SL) systems as additional constraints on top of the standard weak-lensing (WL) shear and
flexion signals.

For the SIS model the combined pipeline had been verified and worked correctly (Task 12).
For the NFW model, however, enabling strong lensing
(\texttt{use\_strong\_lensing=True}) \emph{degraded} the position
reconstruction compared to running the WL-only pipeline.  A concrete test scenario
(Task 14) demonstrated this clearly:

\begin{itemize}
  \item True NFW halo: position $(x,y)=(5'',-3'')$, mass $M = 5\times10^{14}\,M_\odot$,
        $z_l = 0.3$, field half-width $x_\mathrm{max}=120''$, $N_\mathrm{WL}=100$ sources,
        one two-image strong system with $\sigma_\theta = 0.05''$.
  \item WL-only result: nearest recovered halo typically $\Delta \lesssim 5''$ from truth.
  \item WL+SL result (before fixes): nearest recovered halo $\sim 80$--$130''$ from truth.
\end{itemize}

Adding a strong-lensing constraint should by construction only improve or at worst leave
unchanged the position reconstruction; that it \emph{worsened} it by more than an order of
magnitude indicated deep architectural problems.  Investigation identified three independent
but compounding bugs, each described in its own section below.

% ─────────────────────────────────────────────────────────────────
\section{Bug 1 --- Extreme \texorpdfstring{$\lambda_\mathrm{SL}$}{lambda\_SL} at the NFW Initial Guess}
\label{sec:bug1}
% ─────────────────────────────────────────────────────────────────

\subsection{The combined chi-squared objective}

The pipeline minimises
\begin{equation}
  \chi^2_\mathrm{total} = \chi^2_\mathrm{WL} + \lambda_\mathrm{SL}\,\chi^2_\mathrm{SL},
  \label{eq:total_chi2}
\end{equation}
where $\lambda_\mathrm{SL}$ is pre-computed once at the initial guess via
\begin{equation}
  \lambda_\mathrm{SL} = \frac{\tilde\chi^2_\mathrm{WL}}{\tilde\chi^2_\mathrm{SL}}
  = \frac{\chi^2_\mathrm{WL}/\nu_\mathrm{WL}}{\chi^2_\mathrm{SL}/\nu_\mathrm{SL}},
\end{equation}
and held fixed throughout all subsequent optimisation calls.  The intent is to
equalise the reduced chi-squared contributions of the two terms so neither dominates.

\subsection{Why NFW produces extreme \texorpdfstring{$\lambda$}{lambda}}

For SIS the initial guess comes from the flexion-inferred lens position, which is
typically $\mathcal{O}(10'')$--$\mathcal{O}(50'')$ from truth, giving
$\tilde\chi^2_\mathrm{WL}\approx 4600$ and $\tilde\chi^2_\mathrm{SL}\approx 1000$, so
$\lambda \approx 4.4$ --- a balanced weight.

For NFW the initial-guess routine (\texttt{generate\_initial\_guess}) places each lens
candidate at the flexion-predicted lens position \emph{per source}, which is
statistically near the true halo.  At this initial position the strong-lensing images
are nearly back-projected to the same source position, giving
$\tilde\chi^2_\mathrm{SL} \approx 1$--$2$.  Meanwhile the WL chi-squared at the initial
guess --- before mass optimisation --- is large:
$\tilde\chi^2_\mathrm{WL}\approx 500$--$800$.  The ratio is therefore
\begin{equation}
  \lambda_\mathrm{SL}^\mathrm{raw} \approx \frac{700}{1.5} \approx 500,
\end{equation}
a weight two orders of magnitude larger than the SIS case.  With this weight the
entire optimisation effectively minimises $\chi^2_\mathrm{SL}$ only, driving every
halo candidate toward the SL minimum regardless of what the WL data prefer.

\subsection{Fix: cap \texorpdfstring{$\lambda_\mathrm{SL}$}{lambda\_SL} at 50}

In \texttt{arch/metric.py}, function \texttt{compute\_lambda\_sl}:

\begin{lstlisting}
lambda_raw = rchi2_wl / rchi2_sl
lambda_max = 50.0
result = min(lambda_raw, lambda_max)
cap_note = f"  (capped at {lambda_max:.0f})" if lambda_raw > lambda_max else ""
print(f"Pre-computed lambda_sl: {rchi2_wl:.3f} / {rchi2_sl:.3f} = {lambda_raw:.3f}{cap_note}")
return float(result)
\end{lstlisting}

The cap of 50 prevents SL from overwhelming WL when the initial guess happens to
satisfy SL well.  A value of 50 allows SL to be a significant but not dominant
constraint; it was chosen so that the WL reduced chi-squared and the weighted SL
contribution are of comparable magnitude for typical NFW scenarios.

% ─────────────────────────────────────────────────────────────────
\section{Bug 2 --- Non-Monotonic \texorpdfstring{$\chi^2_\mathrm{SL}$}{chi2\_SL} with Magnification Correction}
\label{sec:bug2}
% ─────────────────────────────────────────────────────────────────

\subsection{Source-plane scatter chi-squared with magnification correction}

The strong-lensing chi-squared is computed as the weighted source-plane scatter of
back-projected image positions.  For a system with $n_\mathrm{img}$ images at observed
positions $\boldsymbol\theta_m$, the deflection $\boldsymbol\alpha(\boldsymbol\theta_m)$
gives back-projected positions
\begin{equation}
  \boldsymbol\beta_m = \boldsymbol\theta_m - \boldsymbol\alpha(\boldsymbol\theta_m),
\end{equation}
and the scatter around the weighted mean $\bar{\boldsymbol\beta}$ is
\begin{equation}
  \chi^2_\mathrm{SL} = \sum_m
  \left[\frac{(\beta_{x,m} - \bar\beta_x)^2}{\sigma_{\beta,x,m}^2}
        + \frac{(\beta_{y,m} - \bar\beta_y)^2}{\sigma_{\beta,y,m}^2}\right].
\end{equation}
With the magnification correction enabled, the positional uncertainty in the source
plane is related to the image-plane uncertainty by
\begin{equation}
  \sigma_{\beta,m} = \frac{\sigma_{\theta,m}}{|\mu_m|},
  \qquad |\mu_m| = |\det\mathbf{A}(\boldsymbol\theta_m)|^{-1},
\end{equation}
where $\mathbf{A}$ is the lensing Jacobian.

\subsection{Why this creates a perverse landscape for NFW optimisation}

Near the Einstein radius $\theta_E$ the magnification diverges, $|\mu|\to\infty$,
making $\sigma_\beta \to 0$.  Consequently, \emph{any} mis-deflection, however small,
produces an enormous $\chi^2_\mathrm{SL}$ for lens models whose critical curve passes
near the observed images.

This creates a landscape that is \textbf{non-monotonic in distance from the true lens position}:

\begin{center}
\begin{tabular}{lcc}
  \toprule
  Candidate location & $\chi^2_\mathrm{SL}$ (with mag.\ correction) & Notes\\
  \midrule
  True position, true mass & $\approx 0$ & exact solution\\
  $\sim 6''$ from truth, WL-optimal mass & $\approx 163\,000$ & near critical curve, $|\mu|$ large\\
  $\sim 110''$ from truth & $\approx 2\,100$ & far field, $|\mu|\approx 1$, geometric plateau\\
  \bottomrule
\end{tabular}
\end{center}

The ``geometric plateau'' for two images separated by $2\theta_E \approx 3.3''$ with
$\sigma_\theta = 0.05''$ is
\begin{equation}
  \chi^2_\mathrm{plateau} = \frac{(2\theta_E)^2}{2\sigma_\theta^2}
  = \frac{(3.3'')^2}{2\,(0.05'')^2} \approx 2\,180.
\end{equation}
A halo 110'' from truth therefore has a \emph{lower} $\chi^2_\mathrm{SL}$ than a halo
6'' from truth (but with the wrong mass).  With $\lambda_\mathrm{SL} = 50$, the
contribution $\lambda\chi^2_\mathrm{SL}$ is $50 \times 2\,100 = 105\,000$ for the
far-field plateau but $50 \times 163\,000 = 8\,150\,000$ for the near-truth
candidate.  Consequently:

\begin{enumerate}
  \item \textbf{Forward lens selection} prefers far-from-truth plateau candidates over
        near-truth candidates with even slightly wrong mass.
  \item \textbf{Mass optimisation} (\texttt{optimize\_lens\_strength}) drives the lens
        mass to the minimum allowed value ($10^{10}\,M_\odot$), where the lens barely
        deflects any light and $\chi^2_\mathrm{SL}$ falls back to the plateau value,
        rather than finding the physically correct mass.
\end{enumerate}

\subsection{Fix: use uncorrected chi-squared in all optimisation contexts}

The magnification correction is physically important for the \emph{evaluation} of a
final model (it correctly down-weights poorly-constrained source-plane positions), but
it is harmful as an optimisation objective because it creates the non-monotonic
landscape described above.

The fix introduces a Boolean parameter
\texttt{use\_magnification\_correction\_sl} (default \texttt{True} for backward
compatibility) to \texttt{calculate\_total\_chi2} in \texttt{arch/metric.py}.  All
\emph{optimisation} call sites pass \texttt{use\_magnification\_correction\_sl=False}:

\begin{itemize}[noitemsep]
  \item Every branch of \texttt{chi2wrapper} in \texttt{arch/pipeline.py}
        (called by \texttt{optimize\_lens\_strength} and the SIS position optimiser).
  \item The inner loop of \texttt{forward\_lens\_selection} in
        \texttt{arch/pipeline.py}.
  \item The SL-refinement pass in \texttt{optimize\_lens\_positions} (Steps 5b/5c,
        introduced below).
\end{itemize}

\textit{Evaluation} call sites (\texttt{update\_chi2\_values}, \texttt{compute\_lambda\_sl},
final reporting in \texttt{paper2.py}) retain the default
\texttt{use\_magnification\_correction\_sl=True}.

Without magnification correction the landscape becomes well-behaved:

\begin{center}
\begin{tabular}{lcc}
  \toprule
  Candidate location & $\chi^2_\mathrm{SL}$ (no mag.\ correction) & Notes\\
  \midrule
  True position, true mass & $\approx 0$ & \\
  $1.5''$ from truth, WL-optimal mass & $7.5$ & small scatter\\
  $54''$ from truth & $142$ & growing\\
  $110''$ from truth & $\approx 2\,180$ & geometric plateau\\
  \bottomrule
\end{tabular}
\end{center}

The function is now monotonically increasing with distance from the true position
(approximately), making it a valid optimisation objective.

% ─────────────────────────────────────────────────────────────────
\section{Bug 3 --- Integer Overflow on Infinite Degrees of Freedom}
\label{sec:bug3}
% ─────────────────────────────────────────────────────────────────

\subsection{Origin of the overflow}

\texttt{calc\_degrees\_of\_freedom} returns \texttt{np.inf} when the number of
observational constraints is less than or equal to the number of free parameters ---
e.g.\ when \texttt{filtered\_sources} retains fewer than two WL sources for a
single NFW lens with three free parameters.  The function \texttt{calculate\_total\_chi2}
then attempts to convert this to an integer:

\begin{lstlisting}
# BEFORE (raises OverflowError: cannot convert float infinity to integer)
dof_total = int(dof_wl) + int(dof_sl)
components = {"dof_wl": int(dof_wl), ...}
\end{lstlisting}

\subsection{Fix: guard against infinite DOF}

\begin{lstlisting}
_dof_wl_int = int(dof_wl) if np.isfinite(dof_wl) else 0
_dof_sl_int = int(dof_sl) if np.isfinite(dof_sl) else 0
dof_total = _dof_wl_int + _dof_sl_int

components = {
    "chi2_wl": float(chi2_wl),
    "chi2_sl": float(chi2_sl),
    "dof_wl": _dof_wl_int,
    "dof_sl": _dof_sl_int,
    "lambda_sl": float(_lambda),
}
\end{lstlisting}

Mapping infinite DOF to zero is safe: a zero DOF means the reduced chi-squared is
reported as \texttt{np.inf}, signalling that the model is unconstrained rather than
crashing the pipeline.

% ─────────────────────────────────────────────────────────────────
\section{Residual Failure After Bugs 1--3: Merging Destroys the Near-Truth Halo}
\label{sec:residual}
% ─────────────────────────────────────────────────────────────────

After applying the three fixes above, Tests 11--13 all passed and the SL pipeline no
longer crashed.  However, Task 14-B --- which requires the WL+SL result to be closer
to truth than the WL-only result --- still failed intermittently.  The failure mode
changed: instead of placing the halo 100''+ from truth, the pipeline now produced a
result 20--30'' from truth.

\subsection{Root cause: merge corrupts the near-truth candidate}

Tracing the pipeline step by step revealed the following sequence:

\begin{enumerate}
  \item \textbf{Candidate generation and WL optimisation (Steps 1--2)} produce $\sim 70$
        candidates, with at least one landing within $\sim 1.5''$ of the true halo.
        Each candidate's WL chi-squared with uncorrected SL is:
        \begin{itemize}[noitemsep]
          \item Near-truth candidate: $\chi^2_\mathrm{SL}^\mathrm{nomag} = 7.5$
          \item Next-best: $142$
          \item Far-field plateau: $\approx 2\,200$
        \end{itemize}
        The near-truth candidate is clearly preferred by the combined no-mag objective.
  \item \textbf{Forward selection (Step 4)} correctly selects the near-truth candidate
        first.  However, it then \emph{continues} to add further candidates that each
        reduce $\chi^2_\mathrm{WL}$ marginally (WL overfitting).  Once the near-truth
        halo is in the set, adding a distant wrong halo barely changes $\chi^2_\mathrm{SL}$
        (the near-truth halo already satisfies SL; the wrong halo is a far-field
        perturbation) but does improve $\chi^2_\mathrm{WL}$.  The adaptive tolerance
        is too small to prevent this.
  \item \textbf{First merge (Step 5)} uses a threshold
        \begin{equation}
          d_\mathrm{merge} = \left(\frac{N_\mathrm{src}}{(2x_\mathrm{max})^2}\right)^{-1/2}
          \approx 24''
        \end{equation}
        for $N_\mathrm{src}=100$, $x_\mathrm{max}=120''$.  If any of the wrongly added
        candidates lands within 24'' of the near-truth halo, the two are mass-weighted
        merged.  With the wrong halo potentially having a large WL-optimal mass, the
        merged position can be pulled tens of arcseconds from the true halo, destroying
        the good result.
\end{enumerate}

% ─────────────────────────────────────────────────────────────────
\section{Architectural Fix: Two-Phase Position Refinement (Steps 5b and 5c)}
\label{sec:arch}
% ─────────────────────────────────────────────────────────────────

\subsection{Design principle}

The key insight is that the magnification-corrected $\chi^2_\mathrm{SL}$ is suitable for
\emph{evaluating} a final model but unsuitable as an \emph{optimisation} objective, while
the uncorrected $\chi^2_\mathrm{SL}$ is monotonically well-behaved with a single global
minimum at the true lens position and grows smoothly to the far-field plateau.

Furthermore, the WL chi-squared computed against \emph{all} sources (rather than only the
filtered 20'' neighbourhood used in the initial pass) has a clear global minimum at the
true halo position, because the field signal is dominated by the true halo.

Combining these two properties, the objective
\begin{equation}
  \mathcal{L}_\mathrm{refine}(\boldsymbol{x}, M)
  = \chi^2_\mathrm{WL}(\text{all sources}) + \lambda_\mathrm{SL}\,\chi^2_\mathrm{SL}^\mathrm{nomag}
  \label{eq:refine_obj}
\end{equation}
has a well-defined global minimum at the true lens position with the true mass, and is
suitable for gradient-based optimisation using L-BFGS-B.

\subsection{Implementation}

Two new pipeline steps are inserted in \texttt{arch/main.py} \emph{after} the first merge
and \emph{before} the final mass optimisation.  They only execute when
\texttt{use\_strong\_lensing=True} and \texttt{lens\_type='NFW'}.

\paragraph{Step 5b --- SL-aware position+mass refinement.}
\texttt{optimize\_lens\_positions} is called with \texttt{all\_sources=True}
(no 20'' filter) and \texttt{use\_strong\_lensing=True}.  A new branch in the
objective function detects the combination \texttt{all\_sources=True, use\_strong\_lensing=True}
and uses \eqref{eq:refine_obj} (no mag.\ correction) instead of the WL-only objective
used in the initial pass.  Each selected halo (including any that were merged to a wrong
position) is re-optimised over $(x, y, \log_{10}M)$ simultaneously.  Because both WL and
SL gradients point toward the true halo position, L-BFGS-B converges rapidly toward truth.

\paragraph{Step 5c --- Re-merge.}
After Step 5b, halos that started at different positions may have converged to nearly the
same location.  A second call to \texttt{merge\_close\_lenses} with the same threshold
collapses these into a single halo, preventing double-counting of the mass in Step 6.

\subsection{Why Step 5b does not reintroduce the original Bug 2 failure}

In the \emph{initial} candidate generation (Step 2), $\sim 100$ candidates are optimised
with filtered sources.  If SL were included there, each candidate's SL gradient would pull
it toward the true position, but the filtered WL gradient (computed only from local sources
near the initial guess, not from sources near the true halo) would oppose this, creating an
inconsistent landscape and chaotic convergence.

In Step 5b there are only $\sim 3$--$6$ halos (after forward selection and merging), and
each is optimised with the \emph{full} source catalogue.  The WL signal is dominated by the
true halo, so the WL gradient also points toward truth for all candidates.  Both terms in
\eqref{eq:refine_obj} therefore reinforce each other, and convergence is robust.

% ─────────────────────────────────────────────────────────────────
\section{Summary of All Code Changes}
\label{sec:changes}
% ─────────────────────────────────────────────────────────────────

\begin{description}[style=nextline]

  \item[\texttt{arch/metric.py} --- \texttt{compute\_lambda\_sl}]
    Added hard cap $\lambda_\mathrm{SL} \le 50$ to prevent SL from overwhelming WL when
    the initial guess happens to satisfy the strong constraints well.

  \item[\texttt{arch/metric.py} --- \texttt{calculate\_total\_chi2}]
    \begin{enumerate}[noitemsep]
      \item Added parameter \texttt{use\_magnification\_correction\_sl: bool = True}.
            Passed through to \texttt{chi2\_strong\_source\_plane\_sis/nfw}.
      \item Guarded \texttt{int(dof\_wl)} and \texttt{int(dof\_sl)} against
            \texttt{np.inf} (returns 0 in that case).
    \end{enumerate}

  \item[\texttt{arch/pipeline.py} --- \texttt{optimize\_lens\_positions}]
    \begin{enumerate}[noitemsep]
      \item Added parameter \texttt{all\_sources: bool = False}.
            When \texttt{True}, the 20'' source filter is skipped; all WL sources are used.
      \item Objective function now uses the full WL + no-mag SL objective
            \eqref{eq:refine_obj} when \texttt{all\_sources=True} and
            \texttt{use\_strong\_lensing=True}; otherwise reverts to WL-only
            (preserving the original initial-pass behaviour exactly).
    \end{enumerate}

  \item[\texttt{arch/pipeline.py} --- \texttt{chi2wrapper}]
    All four call sites to \texttt{calculate\_total\_chi2} now pass
    \texttt{use\_magnification\_correction\_sl=False}.

  \item[\texttt{arch/pipeline.py} --- \texttt{forward\_lens\_selection}]
    The inner \texttt{calculate\_total\_chi2} call now passes
    \texttt{use\_magnification\_correction\_sl=False} with a comment explaining why.

  \item[\texttt{arch/main.py} --- \texttt{fit\_lensing\_field}]
    Added Steps 5b and 5c (SL refinement pass and re-merge) after the original Step 5
    merge, conditional on \texttt{use\_strong\_lensing=True} and
    \texttt{lens\_type='NFW'}.

\end{description}

% ─────────────────────────────────────────────────────────────────
\section{Updated Pipeline Flow}
\label{sec:pipeline}
% ─────────────────────────────────────────────────────────────────

Table~\ref{tab:pipeline} summarises the full \texttt{fit\_lensing\_field} pipeline after
the fixes, including which chi-squared objective is used at each step when
\texttt{use\_strong\_lensing=True}.

\begin{table}[h!]
\centering
\caption{Updated \texttt{fit\_lensing\_field} pipeline.  ``Mag.\ corr.'' refers to whether
the magnification correction is applied to $\chi^2_\mathrm{SL}$.  Steps 5b and 5c are new;
all other steps existed previously.}
\label{tab:pipeline}
\small
\begin{tabular}{clccc}
  \toprule
  Step & Action & Sources used & SL included? & Mag.\ corr.?\\
  \midrule
  1 & Generate initial candidates (one per WL source) & --- & No & ---\\
  1b & Pre-compute $\lambda_\mathrm{SL}$ (capped at 50) & All & Yes (evaluation) & Yes\\
  2 & \texttt{optimize\_lens\_positions} (initial pass) & Filtered 20'' & No & ---\\
  3 & \texttt{filter\_lens\_positions} & --- & --- & ---\\
  4 & \texttt{forward\_lens\_selection} & All & Yes (if WL+SL) & \textbf{No}\\
  5 & \texttt{merge\_close\_lenses} (first) & --- & --- & ---\\
  5b & \texttt{optimize\_lens\_positions} (SL refinement, NFW+SL only) & All & Yes & \textbf{No}\\
  5c & \texttt{merge\_close\_lenses} (re-merge, NFW+SL only) & --- & --- & ---\\
  6 & \texttt{optimize\_lens\_strength} (mass) & All & Yes (if WL+SL) & \textbf{No}\\
  7 & \texttt{update\_chi2\_values} (final reporting) & All & Yes (if WL+SL) & Yes\\
  \bottomrule
\end{tabular}
\end{table}

% ─────────────────────────────────────────────────────────────────
\section{Test Results}
\label{sec:tests}
% ─────────────────────────────────────────────────────────────────

All 23 tests in Tasks 11--14 pass after the fixes.  Table~\ref{tab:results} shows
representative results for Task 14 (NFW integration) across multiple test runs.

\begin{table}[h!]
\centering
\caption{Task 14 NFW integration test results (three independent runs with the same
seed).  $\Delta_\mathrm{WL}$ and $\Delta_\mathrm{SL}$ are the distances from the nearest
recovered halo to the true halo position; $M_\mathrm{SL}$ is its mass.
The true values are $\Delta=0$, $M=5\times10^{14}\,M_\odot$.}
\label{tab:results}
\begin{tabular}{ccccc}
  \toprule
  Run & $\Delta_\mathrm{WL}$ ('') & $M_\mathrm{WL}$ ($M_\odot$) & $\Delta_\mathrm{SL}$ ('') & $M_\mathrm{SL}$ ($M_\odot$)\\
  \midrule
  1 & 5.0 & $5.2\times10^{14}$ & \textbf{0.38} & $\mathbf{5.0\times10^{14}}$\\
  2 & 24.7 & $3.2\times10^{14}$ & \textbf{1.36} & $\mathbf{5.9\times10^{14}}$\\
  3 & 1.97 & $5.2\times10^{14}$ & \textbf{0.95} & $\mathbf{5.4\times10^{14}}$\\
  \bottomrule
\end{tabular}
\end{table}

In every run, $\Delta_\mathrm{SL} < \Delta_\mathrm{WL}$: adding strong lensing
consistently improves the position reconstruction, as physically expected.  The recovered
mass with SL is also closer to the true value ($5\times10^{14}\,M_\odot$) in all runs,
reflecting the additional Einstein-radius mass constraint.

Note that the WL-only results themselves vary between runs (e.g.\ $\Delta_\mathrm{WL}$ from
$2''$ to $25''$), reflecting sensitivity of the WL-only pipeline to the specific local
minimum found by L-BFGS-B starting from the flexion-estimated positions.  The SL-augmented
pipeline is more robust: it always finds a halo within $\sim1''$--$2''$ of truth.

% ─────────────────────────────────────────────────────────────────
\section{Remaining Limitations}
\label{sec:limitations}
% ─────────────────────────────────────────────────────────────────

\begin{enumerate}
  \item \textbf{WL overfitting.}  \texttt{forward\_lens\_selection} can still select
        spurious halos because once the true halo is in the set, any additional wrong halo
        barely perturbs $\chi^2_\mathrm{SL}$ (the true halo dominates the deflection) while
        marginally improving $\chi^2_\mathrm{WL}$.  Steps 5b and 5c mitigate the downstream
        damage but do not prevent the spurious candidates from being selected.

  \item \textbf{Non-determinism.}  The WL-only pipeline exhibits run-to-run variability
        (different numbers of halos, positions varying by up to $\sim 20''$), presumably
        from sensitivity of L-BFGS-B to floating-point initialisation order.  The SL
        refinement step reduces this variability but does not eliminate it.

  \item \textbf{SL-only position constraint not yet exploited.}  The current pipeline uses
        SL to \emph{refine} positions found by WL, but does not use SL to \emph{detect}
        halos that the WL pipeline missed entirely (e.g.\ in very noisy fields).

  \item \textbf{Step 5b uses per-halo optimisation.}  Halos are refined individually in
        Step 5b, which is correct when each halo is well separated, but could break down
        for multi-halo systems where the halos contribute comparably to the lensing signal.
\end{enumerate}

\end{document}
